\chapter{Conclusions, Contributions and Implications}
\label{ch:conclusions}

This chapter answers the four research questions directly, states the contributions, sets out
the implications for practice and research, and identifies future work. No new evidence or
references are introduced.

\section{Answers to the research questions}
\label{sec:answers}

\subsection*{RQ1 --- How can discrete-event simulation support monthly workforce-capacity
planning under heterogeneous warehouse demand, and what does it add over an analytical capacity
estimate alone?}

Discrete-event simulation supports the monthly capacity decision by evaluating a bounded set of
workforce candidates against the class-specific service commitments that actually bind, taking
account of queueing, arrival burstiness, stage-differentiated workload and the sequencing rule
in force. The analytical estimate retains a role as a cheap anchor that places the search in the
right region, but it is not the answer.

What simulation adds proved to be more than accuracy. In all twelve months studied, the
simulated recommendation was \emph{smaller} than the analytical estimate --- by up to nine
full-time equivalents in December, where the recommended 38 FTE lay roughly \SI{19}{\percent}
below the estimate of 47. Part of that gap is the deliberate utilisation margin built into the
formula. The remainder arises
because the formula cannot represent a substitution that the simulation discovers: directing
scarce capacity at the orders whose deadlines bind allows a smaller workforce to satisfy the
same commitments. Simulation therefore does not merely compute the same quantity more
carefully; it evaluates a larger decision space in which the operating discipline is a decision
variable rather than a fixed background condition.

\subsection*{RQ2 --- How do alternative sequencing policies affect differentiated service-level
attainment and operating cost?}

The effect is large, and it is conditional on the capacity regime.

Aggregated over 192 runs per policy, FIFO satisfied both service floors in \num{84} while
Urgent-First and RL-3 satisfied them in \num{166} and \num{167}. Its \num{108} infeasible
configurations are spread across all twelve months; what is confined to the four campaign months
is \emph{complete} monthly infeasibility. In January, February, November and
December, FIFO met both floors at \emph{none} of the sixteen candidate configurations, while in
every other month at least eight of the sixteen were feasible. Within the evaluated workforce
range, then, additional capacity did not make FIFO feasible in the campaign months. This is a
statement about the candidate set rather than a law: a sufficiently large workforce would
eventually satisfy both floors under any work-conserving discipline. What the experiment shows
is that the capacity required to do so under FIFO lies beyond the range a planner would
plausibly consider, because an urgent order's queue position bears no relation to its
deadline.

Between the two urgency-aware policies the difference is narrower and regime-dependent. Both
protect the urgent class almost completely; they differ in what that protection costs the
normal class. Where capacity is ample the three policies are practically indistinguishable ---
at October's \rgm{s11\_7\_5} they differed in total cost by \eur{30}. Where the normal-class
floor comes under pressure they separate decisively: at December's recommended configuration,
strict priority drove normal attainment to \SI{73.8}{\percent}, below its floor, and was
infeasible, while the learned policy held it at \SI{85.4}{\percent} at equivalent urgent
attainment and \SI{13.2}{\percent} lower cost.

The systematic comparison over all 192 matched configurations quantifies this regime
dependence rather than illustrating it. Grouped by utilisation at the constraining stage, the
mean cost difference between the two urgency-aware policies rises monotonically from
\eur{537} below \SI{70}{\percent} to \eur{6320} above \SI{95}{\percent} --- a factor of nearly
twelve. The value of the sequencing decision is therefore not a fixed quantity to be estimated
once, but a function of how hard the operation is being run.

The contribution of the learned policy is therefore specific: not better protection of the
urgent class --- which strict priority already achieves, and which RL-3 matched or bettered in
none of the 192 configurations --- but a better outcome for the normal class where protecting
the urgent class becomes counter-productive. That difference is a difference between the two
policies \emph{as implemented}. They differ both in how the serving class is chosen and in how
orders are ordered within a class, and this design cannot decompose the observed gap into those
two effects; the experiment establishes a difference, not a causal estimate of the value of
learning on its own. That advantage is narrow across the year, and it is bounded in a way that
must be reported alongside it: in the 25 configurations where neither policy produced a feasible
plan, the learned policy was worse than strict priority in every one, by nearly 18 percentage
points of normal-class service on average. Its advantage holds where viable plans exist and
reverses where they do not. This study does not establish that either behaviour generalises
beyond the conditions examined.

\subsection*{RQ3 --- How can stage-level bottleneck diagnostics and adaptive capacity evaluation
improve the explainability and economic interpretation of workforce recommendations?}

The diagnostic locates pressure and decomposes the judgement, so that a recommendation arrives
with a reason attached. In this study picking was identified as the primary pressure location in
all twelve months, accounting for between \SI{83.7}{\percent} and \SI{100}{\percent} of the
waiting time accumulated by late orders --- a concentration traceable to the stage-differentiated
workload structure, in which picking carries roughly three times the workload units of dispatch
per order.

The more consequential result concerns what follows from a diagnosis. Pairing the diagnostic
with an explicit economic test showed that identifying a bottleneck does \emph{not} imply that
capacity should be added there. In nine months the framework tested one additional worker at the
identified bottleneck and rejected it in all nine: an FTE costing \eur{2400} per month avoided
less than \eur{130} of penalties in six of them. The recommended configuration has already
absorbed the failures capacity can cheaply prevent, leaving a flat marginal return against a
linear marginal cost.

Explainability is improved on both counts --- the planner learns where the constraint is and,
separately, whether relieving it is worth paying for --- and the second answer is frequently the
opposite of what the first would suggest.

\subsection*{RQ4 --- How can the same simulation framework support both retrospective and
forecast-based planning, and how do their uncertainty structures differ?}

A single engine, feasibility criterion and economic model serve both workflows; they differ only
in how the order stream is obtained. The retrospective workflow replays an order-level history
and answers a counterfactual question, producing a point recommendation per month. The forecast
workflow transforms an aggregate demand assumption into freshly generated, replicated stochastic
order streams and answers a prospective question, producing a spread-aware recommendation ---
mean cost, the cost observed in the least favourable replication, and the share of replications
meeting the service floors --- refined through a screening and validation design that consumed
72 simulation runs against the 144 an exhaustive three-replication evaluation would have needed.
With three replications those figures describe the spread observed rather than estimating a
distribution.

The uncertainty structures are genuinely different and their outputs are not interchangeable.
The retrospective evidence in this study rests on one realisation per month and therefore
carries no replication-based uncertainty estimate; common random numbers make the comparison
between policies fair at a given configuration but say nothing about how the recommendation
would move under different demand. The forecast evidence expresses scenario variation, but rests
on three replications of a single month, so it describes the spread observed rather than
quantifying reliability. That the December forecast run reproduced the retrospective policy pattern under
an independent demand realisation is mild supporting evidence, weakened by the fact that both
runs share the same generator and structural assumptions.

\section{Contributions}
\label{sec:conclusions-contributions}

\textbf{Methodological.} The thesis demonstrates a design in which simulated physical capacity
and economically paid capacity are the same quantity by construction, removing a discrepancy
that would otherwise make every capacity recommendation optimistic; a finite-horizon monthly
formulation that retains unfinished work as explicit backlog instead of allowing unlimited
post-period processing to conceal insufficient capacity; a conjunctive class-specific
feasibility criterion in place of aggregate service attainment; and a three-way policy
comparison under common random numbers in which a learned policy faces exactly the stochastic
realisations that the transparent heuristics face.

\textbf{Practical.} It delivers a workflow that converts either an order-level history or an
aggregate forecast into a recommendation a manager can act on and interrogate: workforce per
stage, sequencing policy, expected cost, class-level attainment, a named pressure location with
its decomposition, and an explicit economic test of the marginal full-time equivalent.

\textbf{Empirical.} It provides evidence that the value of sequencing policy is
capacity-regime dependent --- negligible under slack, decisive under pressure --- and that a
diagnosed bottleneck is routinely not an economically justified place to add capacity. Both
findings run against common intuition and both are directly actionable.

The thesis explicitly does not contribute a new learning algorithm, a globally optimal capacity
solution, or a model validated against industrial data.

\section{Implications}
\label{sec:implications}

\textbf{For practitioners.} Capacity and operating discipline should be sized together, since
the workforce required depends on how work is ordered. Differentiated commitments must be
measured and constrained separately, because aggregate attainment conceals minority-class
failure --- one configuration in this study reported \SI{81}{\percent} overall attainment while
delivering \SI{24.5}{\percent} to urgent customers. Effort spent refining the sequencing rule is
wasted unless the operation runs near its capacity limit, and is essential when it does. A
bottleneck should be treated as a place to investigate rather than a case for hiring. Peak
months should be planned on how consistently a configuration holds across plausible demand, not on
its expected performance.

\textbf{For research.} The result that strict priority can fail a conjunctive class-specific
constraint --- by protecting one class into the failure of the other --- suggests that
evaluating dispatching rules against aggregate objectives such as mean tardiness may understate
the differences that matter when service commitments are contractual and class-specific.
Similarly, the narrowness of the regime in which the learned policy outperformed a one-line
heuristic argues for reporting \emph{where} a learned controller helps rather than whether it
wins on average.

\textbf{For methodology.} The consistent, one-directional gap between the analytical estimate
and the simulated recommendation is a caution against using workload-ratio sizing as anything
more than a starting point when service commitments are differentiated.

\section{Future work}
\label{sec:future-work}

Several directions follow from the limitations of this study, ordered by how directly they
address them.

\textbf{Resolve the within-class ordering asymmetry.} The most valuable single extension would
be an Urgent-First variant maintaining first-in-first-out order within each class, compared
against the current implementation and against RL-3. This would separate the effect of the
priority rule from the effect of within-class ordering, and would establish how much of the
learned policy's advantage survives a fairer baseline.

\textbf{Replicate the retrospective workflow.} Running each month across multiple seeds would
attach uncertainty estimates to the monthly recommendations and establish whether the
near-ties observed in the shoulder months are stable.

\textbf{Calibrate empirically.} Fitting service-time distributions to observed handling data,
and economic parameters to real labour costs and contractual penalties, is the precondition for
any predictive claim.

\textbf{Model shifts explicitly.} Translating monthly FTE into rosters, with shift patterns,
breaks and absence, would close the gap between what the framework recommends and what an
operation can implement.

\textbf{Extend the action space.} The learned policy currently chooses between two classes at
one stage. Richer formulations --- allocating workers across stages, or choosing among more
than two queues --- would test whether the approach scales beyond a binary decision.

\textbf{Sensitivity of the bottleneck diagnostic.} Because one stage dominated in every month
here, the diagnostic's ability to discriminate was never tested. Applying it to a deliberately
balanced operation would establish whether the chosen weights matter.

\textbf{Richer scenario generation and forecast uncertainty.} Propagating genuine forecast error
into the scenario generator, rather than a configured uncertainty level, would make the forecast
workflow's reliability estimates meaningful as predictions rather than as internal consistency
checks.

\section{Closing remarks}
\label{sec:closing}

The question this thesis set out to address was how discrete-event simulation can support the
joint monthly decision of how much labour capacity to deploy and in what order to release work,
under demand that is heterogeneous and service commitments that are differentiated.

The answer that emerges is that the two decisions are not separable, and that treating them as
separable is expensive in a specific and measurable way. Sizing capacity from workload ratios
over-provided by up to nine full-time equivalents in December --- the simulated recommendation
lay approximately nineteen per cent below the analytical estimate --- because the formula could
not represent the fact that a sequencing rule decides which orders wait. Conversely, the value of a sequencing rule proved
almost entirely conditional on how tightly capacity was constrained: irrelevant under slack,
and in the peak season the difference between a viable plan and no viable plan at any workforce
examined.

Two further findings are worth carrying away because they cut against expectation. A stage can
be unambiguously the bottleneck --- saturated, and responsible for almost all the waiting that
late orders accumulate --- and still be the wrong place to add a worker, because the failures
that remain are not the kind that capacity prevents. And a learned policy can be genuinely
valuable in a narrow regime while being indistinguishable, for most of the year, from a rule
that takes one line to state. Reporting where a method helps, rather than whether it wins, seems
to the author the more useful form of both results.

None of this has been validated against a real warehouse, and the thesis is careful not to
claim otherwise. What it offers is an internally coherent instrument for comparing plans under
assumptions that are stated explicitly, together with evidence about the structure of the
problem that a practitioner could test against their own operation.
